\documentclass[UTF8,fontset=fandol,11pt]{ctexart}
\usepackage[a4paper,margin=24mm]{geometry}
\usepackage{amsmath,amssymb,amsthm,graphicx,tabularx,booktabs,hyperref}
\newtheorem{definition}{定义}[section]
\newtheorem{theorem}[definition]{定理}
\newtheorem{example}[definition]{例题}
\newtheorem{exercise}[definition]{练习}
\hypersetup{unicode=true,colorlinks=true,linkcolor=blue}
\title{数列极限入门：定义、唯一性与简单运算}
\date{}
\begin{document}
\maketitle
\section{数列极限的定义与唯一性}\label{sec:S01}
% UNIT S01.U01
\begin{definition}
直观地说，若数列的项随着下标 $n$ 增大而越来越接近某个实数 $A$，就称数列趋近于 $A$。这里的“越来越接近”不是要求从某一项开始单调接近，而是要求误差可以小于任意预先给定的正数。\label{def:sequence-limit}

设 $\{a_n\}$ 是实数数列，$A$ 是实数。如果
\[
\forall\,\varepsilon>0,\quad \exists\,N\in\mathbb{N}^+,\quad \forall\,n\in\mathbb{N}^+,\quad n>N\Longrightarrow |a_n-A|<\varepsilon,
\]
则称数列 $\{a_n\}$ 收敛于 $A$，称 $A$ 为该数列的极限，记作
\[
\lim_{n\to\infty}a_n=A.
\]
其中，$\varepsilon$ 是任意给定的正误差容限，$N$ 可以依赖于 $\varepsilon$，但选定后必须对所有满足 $n>N$ 的正整数 $n$ 同时成立。
\end{definition}
% UNIT S01.U02
\begin{theorem}
数列极限具有唯一性：如果实数数列 $\{a_n\}$ 满足
\[
\lim_{n\to\infty}a_n=A\qquad\text{且}\qquad \lim_{n\to\infty}a_n=B,
\]
其中 $A,B\in\mathbb{R}$，那么
\[
A=B.
\]
\label{thm:limit-uniqueness}
\end{theorem}
% UNIT S01.U03
\begin{proof}
反设 $A\ne B$。取
\[
\varepsilon=\frac{|A-B|}{3}>0.
\]
由 $\lim_{n\to\infty}a_n=A$ 的定义，存在正整数 $N_A$，使得当 $n>N_A$ 时，$|a_n-A|<\varepsilon$。同理，存在正整数 $N_B$，使得当 $n>N_B$ 时，$|a_n-B|<\varepsilon$。

取正整数 $N=\max\{N_A,N_B\}$。当 $n>N$ 时，由三角不等式可得
\[
|A-B|=|(A-a_n)+(a_n-B)|
\leq |A-a_n|+|a_n-B|
<2\varepsilon=\frac{2|A-B|}{3}.
\]
这与 $|A-B|>0$ 矛盾。因此不能有 $A\ne B$，所以 $A=B$。
\end{proof}
% UNIT S01.U04
\begin{example}
用数列极限的定义证明：当 $a_n=\dfrac{1}{n}$，$n\in\mathbb{N}^+$ 时，数列 $\{a_n\}$ 的极限为 $0$。
\end{example}
% UNIT S01.U05
\par\noindent\textbf{解：}任取 $\varepsilon>0$。为了使
\[
\left|a_n-0\right|=\frac{1}{n}<\varepsilon,
\]
只需令 $n>1/\varepsilon$。取正整数 $N$ 满足
\[
N>\frac{1}{\varepsilon}.
\]
这样的 $N$ 对每个给定的 $\varepsilon>0$ 都存在。

当 $n>N$ 时，由 $N>1/\varepsilon$ 可知 $n>1/\varepsilon$，于是
\[
\left|a_n-0\right|=\left|\frac{1}{n}\right|=\frac{1}{n}<\varepsilon.
\]
因此，按照定义，数列 $\{1/n\}$ 收敛于 $0$，即
\[
\lim_{n\to\infty}\frac{1}{n}=0.
\label{eq:reciprocal-limit}
\]
这里 $N$ 可以随着 $\varepsilon$ 的变化而重新选择，但在选定某个 $\varepsilon$ 后，$N$ 不能随着待验证的 $n$ 改变。\par
% UNIT S01.U06
\begin{exercise}
用数列极限的定义证明：当 $a_n=2+\dfrac{3}{n}$，$n\in\mathbb{N}^+$ 时，数列 $\{a_n\}$ 的极限为 $2$。
\end{exercise}
\section*{练习答案}
% UNIT S01.U07
\par\noindent\textbf{解：}任取 $\varepsilon>0$。有
\[
\left|a_n-2\right|=\left|2+\frac{3}{n}-2\right|=\frac{3}{n}.
\]
为了使 $\frac{3}{n}<\varepsilon$，只需令 $n>3/\varepsilon$。取正整数 $N$ 满足
\[
N>\frac{3}{\varepsilon}.
\]
当 $n>N$ 时，有 $n>3/\varepsilon$，从而
\[
\left|a_n-2\right|=\frac{3}{n}<\varepsilon.
\]
因此，根据数列极限的定义，
\[
\lim_{n\to\infty}\left(2+\frac{3}{n}\right)=2.
\]\par
\end{document}
